\documentclass[10pt]{article}
\usepackage[diary]{aditya}

\title{Diary entry for Spectral-Sequences}
\author{Aditya Pahuja}
\date{\today}

\begin{document}
\maketitle
\tableofcontents
\pagebreak

\section{Theory}

\section{Examples and applications}

\subsection{Some results in homological algebra}
\begin{prop}
    [Four lemmas]
    \label{prop:four-lemmas}
    Consider the following commutative diagram
    of modules (or just in some abelian category),
    in which the two rows are exact:
    \begin{center}
        \begin{tikzcd}
	    A \arrow[d,"\alpha"]
	    & B\arrow[l]\arrow[d,"\beta"] 
	    & C\arrow[l]\arrow[d,"\gamma"]
	    & D\arrow[l]\arrow[d,"\delta"] \\
	    A' & B'\arrow[l] & C'\arrow[l] & D'\arrow[l]
        \end{tikzcd}
    \end{center}
    \begin{enumerate}[(a)]
        \ii If $\alpha$ and $\gamma$ are monomorphisms
	and $\delta$ is a epimorphism, then $\beta$ is an monomorphism.
	\ii If $\beta$ and $\delta$ are epimorphisms
	and $\alpha$ is an monomorphism, then $\gamma$ is a epimorphism.
    \end{enumerate}
    (The diagram is drawn backwards to emphasize
    how the double complex is oriented in the plane.)
\end{prop}

\begin{proof}
    By extending via zeroes, we can view the given diagram
    as a first quadrant double complex $E$ which is concentrated in
    finitely many points of the grid,
    so the corresponding (starting-at-horizontal or starting-at-vertical)
    spectral sequences converge to the homology of its total complex.
    If we take horizontal homology first to get to $E^1$,
    then the $E^1$ (hence $E^\infty$) page is all zeroes,
    so the spectral sequence must converge to zero.

    Now we consider the two parts separately and
    pass to $E^1$ by taking vertical homology.

    (a) The $E^1$ page is
    \begin{center}
        \begin{tikzcd}
	    0
	    & \ker\beta \arrow[l,"d^1"']
	    & 0\arrow[l,"d^1"']
	    & \ker\delta \\
	    \coker\alpha
	    & \coker\beta
	    & \coker\gamma & 0\arrow[llu,"d^2"']
        \end{tikzcd}
    \end{center}
    From here on out, the differentials into and out of $\ker\beta$
    are all zeroes: in $E^1$, the differentials
    are the adjacent zero maps; in $E^2$, the incoming differential
    is from $\coker\delta = 0$ and the outgoing one
    goes outside the first quadrant; and in the remaining pages
    the differentials come from or end at a place outside the first quadrant.
    This means $\ker\beta$ persists in $E^\infty$,
    which we have already discovered consists of all zeroes,
    so $\ker\beta = 0$.

    (b) Similarly, the $E^1$ page in this case is
    \begin{center}
        \begin{tikzcd}
	    0 
	    & \ker\beta 
	    & \ker\gamma
	    & \ker\delta \\
	    \coker\alpha
	    & 0
	    & \coker\gamma\arrow[llu,"d^2"']\arrow[l,"d^1"']
	    & 0\arrow[l,"d^1"']
        \end{tikzcd}
    \end{center}
    and the differentials involving $\coker\gamma$ from here on out
    are all zeroes, meaning $\coker\gamma$ persists to $E^\infty$,
    hence is zero.
\end{proof}

\begin{exercise}
    \label{exer:five-lemma}
    Prove the five lemma using a similar spectral sequence argument.
    Recall that the five lemma says that for a commutative diagram
    \begin{center}
        \begin{tikzcd}
	    A \arrow[d,"\alpha"]
	    & B\arrow[l]\arrow[d,"\beta"] 
	    & C\arrow[l]\arrow[d,"\gamma"]
	    & D\arrow[l]\arrow[d,"\delta"]
	    & E\arrow[l]\arrow[d,"\delta"] \\
	    A' & B'\arrow[l] & C'\arrow[l] & D'\arrow[l] & E'\arrow[l]
        \end{tikzcd}
    \end{center}
    in which the two rows are exact,
    $\beta$ and $\delta$ are isomorphisms,
    $\delta$ is an epimorphism, and $\alpha$ is a monomorphism,
    we have that $\gamma$ is an isomorphism.
    (In real life this often just manifests as
    the outer four maps being isomorphisms
    implying that $\gamma$ is an isomorphism.)
\end{exercise}

% >>> apparently circular ?
\begin{prop}
    [Snake lemma]
    \label{prop:snake-lemma}
    Consider the commutative diagram of modules
    \begin{center}
        \begin{tikzcd}
	    0
	    & C\arrow[l]\arrow[d,"\gamma"] 
	    & B\arrow[l,"g"']\arrow[d,"\beta"]
	    & A\arrow[l,"f"']\arrow[d,"\alpha"]
	    & \\
	    & C' & B'\arrow[l,"g'"'] & A'\arrow[l,"f'"'] & 0\arrow[l]
        \end{tikzcd}
    \end{center}
    in which the rows are exact.
    There is an exact sequence
    \begin{equation*}
        \ker f\to \ker\alpha\to\ker\beta\to\ker\gamma
	\overset\partial\too\coker\alpha\to\coker\beta\to\coker\gamma
	\to\coker g'.
    \end{equation*}
\end{prop}

\begin{proof}
    As before, complete the diagram with zeroes
    to make it into a first quadrant double complex $E$.
    Taking horizontal homology gives an $E^1$ page of
    \begin{equation*}
        \begin{matrix}
	    0 & 0 & \ker f\\
	    \coker g' & 0 & 0
        \end{matrix}
    \end{equation*}
    after which the spectral sequence stabilizes,
    so this is equal to the $E^\infty$ page.

    Taking vertical homology instead gives an $E^1$ page of
    \begin{equation*}
        \begin{matrix}
	    \ker\gamma & \ker\beta & \ker\alpha\\
	    \coker\gamma & \coker\beta & \coker\alpha
        \end{matrix}
    \end{equation*}
    Passing to $E^2$ by taking horizontal homology
    must kill $\ker\beta$ and $\coker\beta$
    because in future pages they are not involved in nonzero differentials.
    For a similar reason, the homology at $\ker\alpha$ must be $\ker f$
    and the homology at $\coker\gamma$ must be $\coker g'$,
    since the objects at those locations
    also persist from $E^2$ to $E^\infty$.
    This gives the entire exact sequence
    except for the connecting homomorphism, which we consider now.
    In the $E^2$ page the only nontrivial differential is a map
    \begin{equation*}
	d^2 \colon \ker d^1_{2,0} \to \coker d^1_{1,1},
    \end{equation*}
    which must be an isomorphism because the objects at these locations
    must vanish in the $E^\infty$ page and they will not be
    involved in nontrivial differentials in the future.
    Define the connecting map $\partial$ as the composition
    \begin{equation*}
	\ker\gamma\onto\coker d^1_{1,1}
	\overset\sim\too\ker d^1_{2,0}\into\coker\alpha,
    \end{equation*}
    inverting $d^2$ for the middle map.
    The kernel of $\partial$ is that of the first map in the composition,
    namely the image of the map $d^1_{1,1}\colon\ker\beta\to\ker\gamma$
    in the exact sequence, so
    $\ker\beta\to\ker\gamma\overset\partial\to\coker\alpha$
    is exact at $\ker\gamma$.
    Similarly, the image of $\partial$ is isomorphic to the kernel of the map
    in the exact sequence $d^1_{2,0}\colon\coker\beta\to\coker\beta$,
    proving exactness at $\coker\alpha$.
\end{proof}

\begin{exercise}
    \label{exer:homology-LES}
    Use essentially the same argument to show that
    a short exact sequence of chain complexes
    induces a long exact sequence in homology.
\end{exercise}

Note that in all the examples above,
the total complex itself is essentially irrelevant;
the only important thing is that
both spectral sequences furnished by the double complex
converge to the same thing (which fortunately was all zeroes).

\subsection{Serre spectral sequence}

\begin{theorem}[Serre spectral sequence]
    \label{thm:serre-ss}
    Let $F\into E\to B$ be a fibration with path connected base space $B$
    whose fundamental group acts trivially on $H_*(F)$.
    Then, there is a homological spectral sequence
    $E^2_{p,q} = H_p(B; H_q(F))$ converging to $H_*(E)$
    and a cohomological spectral sequence
    $E_2^{p,q} = H^p(B; H^q(F; R))$ converging to $H^*(E)$
    for some abelian group $R$.
    Furthermore, if $R$ is a ring, there are
    bilinear products $E^{p,q}_r\tensor E^{s,t}_r\to E^{p+q,s+t}_r$
    on each page such that
    \begin{enumerate}[(a)]
        \ii Each $d_r$ is a derivation, namely
	$d_r(xy) = (d_rx)y + (-1)^{p+q}xd_ry$ where $x\in E_r^{p+q}$.
	Note that this means the product on page $r+1$
	is induced by the product on page $r$.
	\ii The product in $E^2$ is the standard cup product
	with coefficients in the cohomology ring of the fiber.
	\ii The cup product on $H^*(E)$
	induces the products on the $E^\infty$ page.
    \end{enumerate}
\end{theorem}

This is particularly useful when you can say something nice
about some of the terms of the spectral sequence,
such as knowing that $E$ is contractible
or that there are a lot of zero rows/columns
(perhaps the spectral sequence is concentrated
in just a few rows or columns).
Note also that the triviality of the action of $\pi_1$
often manifests as $X$ just being simply connected.

Here are a whole bunch of examples.

\begin{theorem}[Hurewicz's theorem]
    \label{thm:hurewicz}
    Suppose that $X$ is $n$-connected,
    meaning the first $n$ homotopy groups of $X$ are zero.
    Then, the first $n$ reduced homology groups (with integer coefficients)
    of $X$ are also zero, and
    if $n > 0$, then $\pi_{n+1}(X,x_0)\cong H_{n+1}(X)$.
    If $n = 0$, we instead have $\pi_1(X,x_0)^{\mathrm{ab}}\cong H_1(X)$,
    the abelianization of the fundamental group.
\end{theorem}

\begin{proof}
    (This proof was shown to me by my friend Nick.)

    For $n = 0$ one can prove the theorem manually without
    too much blood, sweat, or tears;
    see Hatcher appendix 2A. To extend to larger $n$, we induct on $n$.
    Suppose that $X$ is $n$-connected for $n\ge 1$
    (notably, simply connected), and consider
    the Serre spectral sequence for the path fibration
    $\Omega X\to PX\to X$. Using the inductive hypothesis
    and the fact that $\Omega X$ is $(n-1)$-connected,
    the $E^2$ page of the Serre spectral sequence looks like
    \begin{equation*}
        \begin{matrix}
            H_n(\Omega X) = \pi_n(\Omega X)\\
	    0 & 0 & \cdots & 0 & 0\\
	    \vdots & \vdots & \ddots & \vdots & \vdots\\
	    0 & 0 & \cdots & 0 & 0\\
	    \ZZ & 0 &\cdots& 0 & H_n(X) & H_{n+1}(X)
        \end{matrix}
    \end{equation*}
    where rows $1$ through $n-1$ are all zeroes
    because the coefficient group $H_k(\Omega X)$ is zero for $1\le k < n$.
    Since the entire diagonal $p+q=n-2$ consists of zeroes,
    $H_n(X)$ can only ever map into zero,
    so it persists in the $E^\infty$ page;
    since $PX$ is contractible, this means $H_n(X)$ must be zero.
    To deal with $H_{n+1}(X)$, we note that
    the entry above $H_n(X)$ is zero by hypothesis if $n > 1$,
    and if $n = 1$ then this entry sits directly to the right of $\pi_n(X)$,
    so there is no differential from it into $\pi_n(X)$.
    Thus the only nontrivial differential involving
    $H_{n+1}(X)$ or $\pi_n(X)$ is the one in page $n+1$,
    where $H_{n+1}(X)$ maps into $\pi_n(X)$;
    this is the only opportunity for these two to become zero
    (which they must in $E^\infty$), which means that
    \begin{equation*}
	0\to H_{n+1}(X) \overset{d^{n+1}}\too \pi_n(\Omega X)\to 0
    \end{equation*}
    is exact, i.e. $H_{n+1}(X)\cong \pi_n(\Omega X) \cong \pi_{n+1}(X)$.
\end{proof}

\begin{example}
    \label{ex:KG1-exact-seq}
    Let $1\to A\to B\to C\to 1$ be a short exact sequence of groups.
    Each homomorphism $B\to C$ can be realized by
    a map $K(B,1)\to K(C,1)$. Converting this into a fibration\todo{elaborate}
    and feeding it into the long exact sequence shows that
    \begin{equation*}
        1\to \pi_1(F)\to B\to C\to 1
    \end{equation*}
    is exact, so $\pi_1(F) = A$.
    \todo{uh ok this is more complicated than i thought}
\end{example}



\section{Practice}
Courtesy of Ajay.

\begin{exercise}
    \label{exer:K-th-CPinfty}
    Use the Atiyah-Hirzebruch spectral sequence
    to compute the $K$-theory of $\CP^\infty$.
\end{exercise}

\begin{exercise}
    \label{exer:CPinfty-homology}
    Compute the homology of $K(\ZZ,2) = \CP^\infty$
    using the Serre spectral sequence,
    noting that there is a fibration $S^\infty\to\CP^\infty$
    (what is its fiber?).
    Note that you should be able to verify your answer easily
    using cellular homology.
\end{exercise}

\begin{exercise}
    \label{exer:RPinfty-homology}
    Do the same as the above exercise for $\RP^\infty$
    (perhaps with $\ZZ_2$ coefficients;
    I am not sure how tractible this is).
\end{exercise}

\begin{exercise}
    \label{exer:covering-space-homology}
    If $E\to B$ is a covering space,
    how are the homology of $E$ and $B$ related?
\end{exercise}

\begin{exercise}
    \label{exer:kunneth}
    Deduce the K\"unneth formula
    for (co)homology with coefficients in a field
    using the Serre spectral sequence.

    Bonus question: What about the case where the
    coefficient ring is not a field?
\end{exercise}

\begin{exercise}
    \label{exer:loop-sphere}
    Compute the cohomology of $\Omega S^n$,
    the loop space of the $n$-sphere.
    What is its cup product structure?
\end{exercise}

\begin{exercise}
    \label{exer:contractible-fiber}
    If $E\to B$ is a fibration whose fiber is contractible,
    what can be said about the homology of $E$ and $B$?
    (For a spoiler of the answer, feed the fibration
    into the homotopy long exact sequence.)
    What about if $B$ is contractible, or $E$ is contractible?
\end{exercise}

Some more, from Nick:

\begin{exercise}
    \label{exer:universal-coefs}
    Prove the various universal coefficient theorems.
\end{exercise}

\begin{exercise}
    \label{exer:ext-tor-well-defined}
    Show that $\ext$ and $\tor$ are well-defined.
    That is, the two right derived functors
    which define $\ext_R^n(A,B)$ agree
    and the two left derived functors
    which define $\tor^R_n(A,B)$ agree.
\end{exercise}

\end{document}

